\documentclass[10pt,a4paper]{amsart}
\usepackage{amsmath,amssymb,amsthm,graphicx}
% Retained minimal preamble for the Stacks Project chapter more-algebra.tex (More on Algebra).
% The chapter file pinned-original-source.tex carries only \input{preamble} on its first line, so no
% preamble text lives inside the chapter itself.  The Stacks Project uses one single shared preamble for
% every chapter; the snapshot stages it identically as sources/stacks-chow/preamble.tex and
% sources/stacks-etale-cohomology/preamble.tex.  The block below is copied mechanically, line for line,
% from sources/stacks-chow/preamble.tex lines 14-18 (the Xy-pic setup) and 154-217 (the native theorem
% declarations together with the mathematical macros).  Nothing was added, removed or redefined; the
% wrapper class and the amsmath/amssymb/amsthm/graphicx base packages are supplied by the assembler.
% Chapter: more-algebra.tex -- More on Algebra.
% For commutative diagrams we use Xy-pic
\usepackage[all]{xy}

% We use 2cell for 2-commutative diagrams.
\xyoption{2cell}

% Theorem environments.
%
\theoremstyle{plain}
\newtheorem{theorem}[subsection]{Theorem}
\newtheorem{proposition}[subsection]{Proposition}
\newtheorem{lemma}[subsection]{Lemma}

\theoremstyle{definition}
\newtheorem{definition}[subsection]{Definition}
\newtheorem{example}[subsection]{Example}
\newtheorem{exercise}[subsection]{Exercise}
\newtheorem{situation}[subsection]{Situation}

\theoremstyle{remark}
\newtheorem{remark}[subsection]{Remark}
\newtheorem{remarks}[subsection]{Remarks}

\numberwithin{equation}{subsection}

% Macros
%
\def\lim{\mathop{\mathrm{lim}}\nolimits}
\def\colim{\mathop{\mathrm{colim}}\nolimits}
\def\Spec{\mathop{\mathrm{Spec}}}
\def\Hom{\mathop{\mathrm{Hom}}\nolimits}
\def\Ext{\mathop{\mathrm{Ext}}\nolimits}
\def\SheafHom{\mathop{\mathcal{H}\!\mathit{om}}\nolimits}
\def\SheafExt{\mathop{\mathcal{E}\!\mathit{xt}}\nolimits}
\def\Sch{\mathit{Sch}}
\def\Mor{\mathop{\mathrm{Mor}}\nolimits}
\def\Ob{\mathop{\mathrm{Ob}}\nolimits}
\def\Sh{\mathop{\mathit{Sh}}\nolimits}
\def\NL{\mathop{N\!L}\nolimits}
\def\CH{\mathop{\mathrm{CH}}\nolimits}
\def\proetale{{pro\text{-}\acute{e}tale}}
\def\etale{{\acute{e}tale}}
\def\QCoh{\mathit{QCoh}}
\def\Ker{\mathop{\mathrm{Ker}}}
\def\Im{\mathop{\mathrm{Im}}}
\def\Coker{\mathop{\mathrm{Coker}}}
\def\Coim{\mathop{\mathrm{Coim}}}

% Boxtimes
%
\DeclareMathSymbol{\boxtimes}{\mathbin}{AMSa}{"02}

%
% Macros for moduli stacks/spaces
%
\def\QCohstack{\mathcal{QC}\!\mathit{oh}}
\def\Cohstack{\mathcal{C}\!\mathit{oh}}
\def\Spacesstack{\mathcal{S}\!\mathit{paces}}
\def\Quotfunctor{\mathrm{Quot}}
\def\Hilbfunctor{\mathrm{Hilb}}
\def\Curvesstack{\mathcal{C}\!\mathit{urves}}
\def\Polarizedstack{\mathcal{P}\!\mathit{olarized}}
\def\Complexesstack{\mathcal{C}\!\mathit{omplexes}}
% \Pic is the operator that assigns to X its picard group, usage \Pic(X)
% \Picardstack_{X/B} denotes the Picard stack of X over B
% \Picardfunctor_{X/B} denotes the Picard functor of X over B
\def\Pic{\mathop{\mathrm{Pic}}\nolimits}
\def\Picardstack{\mathcal{P}\!\mathit{ic}}
\def\Picardfunctor{\mathrm{Pic}}
\def\Deformationcategory{\mathcal{D}\!\mathit{ef}}

% --- appended 2026-10-12 (Chief): the project own definitions of the reference/slogan/history
% display environments, copied mechanically from the shared Stacks preamble lines 8-12
% (\usepackage{verbatim} plus the three \newenvironment lines, which the chapter preamble keeps
% verbatim from the project preamble). Nothing added, removed or redefined: these environments
% carry no mathematical content and the project itself typesets them as comments. Without these
% three lines a counted statement whose own text uses them fails with Environment slogan undefined.
\usepackage{verbatim}
\newenvironment{reference}{\comment}{\endcomment}
\newenvironment{slogan}{\comment}{\endcomment}
\newenvironment{history}{\comment}{\endcomment}
\usepackage{cleveref}
\crefname{originallemma}{original lemma}{original lemmas}
\Crefname{originallemma}{Original lemma}{Original lemmas}
\makeatletter
\newlabel{lemma-derived-completion-koszul}{{0920}{source}{}{sourcepointer-1}{}}
\newlabel{lemma-derived-completion-koszul@cref}{{[originallemma][0][]0920}{[1][0][]source}{}{}{}}
\makeatother
\crefname{originallemma}{original lemma}{original lemmas}
\Crefname{originallemma}{Original lemma}{Original lemmas}
\makeatletter
\newlabel{lemma-break-long-exact-sequence-modules}{{0CQE}{source}{}{sourcepointer-2}{}}
\newlabel{lemma-break-long-exact-sequence-modules@cref}{{[originallemma][0][]0CQE}{[1][0][]source}{}{}{}}
\makeatother
\title{Derived completion has finite cohomological dimension bounded by the number of generators of the ideal}
\author{The Stacks Project Authors}
\date{stacks-project:more-algebra.tex:local-snapshot-20261012; GNU FDL 1.2 or later}
\begin{document}
\maketitle
\setcounter{subsection}{1}
\noindent\textit{Editorial note.} All delivered mathematical text is the author's own, copied from the pinned Stacks Project chapter file more-algebra.tex; no statement or proof step was rewritten, reordered or removed. The counted statement is the labeled lemma environment at source lines 26507-26522 (canonical tag 0AAJ) and the counted proof is the proof environment at source lines 26524-26553, which immediately follows it and is its own proof.\par\medskip
\noindent\textit{Editorial note.} The preamble span runs to line 90, the current end of the retained preamble, which carries the shared Stacks preamble together with the appended environment definitions; the counted text does not use those environments, but the span is taken whole so that the file version is unambiguous.\par\medskip
\noindent\textit{Editorial note.} The printed text of part (2) concludes that the cohomology of the completion of the cone vanishes below degree r while the statement predicts the shift minus r; the printed text is copied exactly as it stands and no correction is applied, since the lane delivers the author's own words.\par\medskip
\noindent\textit{Editorial note.} The proof refers twice to earlier results of this chapter, namely the Koszul description of the completion and the exact sequence for the cohomology of a derived limit; both are navigation stubs with their roles recorded in the reasons.\par\medskip
\noindent\textit{Editorial note.} The source URL is the source-level Stacks Project URL, as required when more than one card is produced from one source; the per-claim tag is carried by claim\_key and by this note. The counted statement and proof cite no bibliography entry, so this card needs no bibliography.\par\medskip
\noindent\textit{Allowed standard prerequisites (editorial).}\begin{itemize}
\item Derived completion of an object is the derived limit of its derived tensor products with the Koszul complexes on the powers of a generating set. invoked at proof lines 26525-26528. Turns the completion into a derived limit so that its cohomology can be computed from the Koszul system.
\item For a system of complexes there is a short exact sequence expressing the cohomology of the derived limit through the first derived limit and the limit of the cohomologies. invoked at proof lines 26528-26532. Provides the displayed short exact sequence from which the vanishing of the cohomology of the completion is read off.
\item The Koszul complex on the n-th powers of r generators sits in degrees from -r to 0, so tensor products with it vanish in the corresponding range of degrees. used at proof lines 26537-26541 and 26547-26550. Gives the degree ranges in which the cone of the given morphism has vanishing cohomology after tensoring with the Koszul system, which pins down the vanishing range for the completion.
\end{itemize}\medskip
\ifdefined\hypertarget\hypertarget{sourcepointer-1}{}\fi
\noindent\textit{Source reference.} Original lemma 0920: sources/stacks-more-algebra/pinned-original-source.tex, source lines 26388--26388. The reference gives the description of the completion as a derived limit of Koszul complexes, used at proof lines 26525-26527; the exact sequence and the two case analyses that follow are inside the counted span.\par\medskip
\ifdefined\hypertarget\hypertarget{sourcepointer-2}{}\fi
\noindent\textit{Source reference.} Original lemma 0CQE: sources/stacks-more-algebra/pinned-original-source.tex, source lines 23653--23653. The reference supplies the displayed short exact sequence for the cohomology of the completion, invoked at proof lines 26528-26532; everything read off from that sequence is inside the counted span.\par\medskip
\begin{lemma}
\label{lemma-derived-completion-finite-cohomological-dimension}
Let $A$ be a ring and let $I \subset A$ be an ideal which can be
generated by $r$ elements. Then derived completion has finite
cohomological dimension:
\begin{enumerate}
\item Let $K \to L$ be a morphism in $D(A)$ such that $H^i(K) \to H^i(L)$
is an isomorphism for $i \geq 1$ and surjective for $i = 0$.
Then $H^i(K^\wedge) \to H^i(L^\wedge)$ is an isomorphism for $i \geq 1$
and surjective for $i = 0$.
\item Let $K \to L$ be a morphism of $D(A)$ such that $H^i(K) \to H^i(L)$
is an isomorphism for $i \leq -1$ and injective for $i = 0$.
Then $H^i(K^\wedge) \to H^i(L^\wedge)$ is an isomorphism for $i \leq -r - 1$
and injective for $i = -r$.
\end{enumerate}
\end{lemma}

\begin{proof}
Say $I$ is generated by $f_1, \ldots, f_r$. For any $K \in D(A)$
by Lemma \ref{lemma-derived-completion-koszul} we have
$K^\wedge = R\lim K \otimes_A^\mathbf{L} K_n$ where $K_n$
is the Koszul complex on $f_1^n, \ldots, f_r^n$
and hence we obtain a short exact sequence
$$
0 \to R^1\lim H^{i - 1}(K \otimes_A^\mathbf{L} K_n)
\to H^i(K^\wedge) \to \lim H^i(K \otimes_A^\mathbf{L} K_n) \to 0
$$
by Lemma \ref{lemma-break-long-exact-sequence-modules}.

\medskip\noindent
Proof of (1). Pick a distinguished triangle $K \to L \to C \to K[1]$.
Then $H^i(C) = 0$ for $i \geq 0$. Since $K_n$ is sitting in degrees $\leq 0$
we see that $H^i(C \otimes_A^\mathbf{L} K_n) = 0$ for $i \geq 0$
and that $H^{-1}(C \otimes_A^\mathbf{L} K_n) =
H^{-1}(C) \otimes_A A/(f_1^n, \ldots, f_r^n)$ is a system with
surjective transition maps. The displayed equation above shows
that $H^i(C^\wedge) = 0$ for $i \geq 0$. Applying the distinguished triangle
$K^\wedge \to L^\wedge \to C^\wedge \to K^\wedge[1]$ we get (1).

\medskip\noindent
Proof of (2). Pick a distinguished triangle $K \to L \to C \to K[1]$.
Then $H^i(C) = 0$ for $i < 0$. Since $K_n$ is sitting in degrees
$-r, \ldots, 0$ we see that $H^i(C \otimes_A^\mathbf{L} K_n) = 0$
for $i < -r$. The displayed equation above shows
that $H^i(C^\wedge) = 0$ for $i < r$. Applying the distinguished triangle
$K^\wedge \to L^\wedge \to C^\wedge \to K^\wedge[1]$ we get (2).
\end{proof}

\end{document}
